% Generated by report/tools/tables.py from the result files of the repository. Do not edit.
\begin{table}[tbp]
\centering
\footnotesize
\setlength{\tabcolsep}{3.5pt}
\caption[Trajectory 2018 to 2026]{Trajectory of classical RL from 2018 to 2026 (deviation 25 of \file{study/PROTOCOL.md}, written before any draw): share of papers that meet each criterion, in percent, by year. 2018: our reading of the papers with experiments; 2019, 2020, 2022 and 2024: the first 20 experimental classical RL papers per year, drawn at random from ICML and NeurIPS (ICLR is missing, its site refusing automated downloads); 2026: the classical sample of the main comparison. The number in parentheses gives the denominator when it differs from the column size (criterion not applicable to some papers). With 20 papers, a share carries an uncertainty of about twenty points either way. For the five tested criteria (Q5, Y9, X7, X2, Q8): Cochran-Armitage trend test over the six years (scores equal to the years), two-sided $p$ and Holm-corrected $p$ over the five; the six other rows are descriptive. Bottom part: immediate gap (2019 and 2020 pooled against 2018) and later gap (2022 and 2024 against 2019 and 2020), in points, 95\% Newcombe interval, Holm over the five criteria within each family. \textbf{In bold}, what passes the Holm correction. Source: \file{study/trajectory.json} (\file{study/trajectory.py}).}
\label{tab:trajectoire}
\begin{tabularx}{\linewidth}{@{}l>{\raggedright\arraybackslash}Xrrrrrrrr@{}}
\toprule
Code & Criterion & 2018 & 2019 & 2020 & 2022 & 2024 & 2026 & trend $p$ & Holm $p$ \\
 & $n$ & 48 & 20 & 20 & 20 & 20 & 50 & &  \\
\midrule
\textbf{Q5} & \textbf{Code linked} & \textbf{31} & \textbf{60} & \textbf{50} & \textbf{70} & \textbf{75} & \textbf{74} & \textbf{2\,×\,10\textsuperscript{−5}} & \textbf{7\,×\,10\textsuperscript{−5}} \\
\textbf{Y9} & \textbf{Sensitivity to a hyperparameter} & \textbf{44} & \textbf{35} & \textbf{45} & \textbf{60} & \textbf{85} & \textbf{80} & \textbf{2\,×\,10\textsuperscript{−6}} & \textbf{9\,×\,10\textsuperscript{−6}} \\
\textbf{X7} & \textbf{Compute quantified} & \textbf{27} & \textbf{30} & \textbf{40} & \textbf{40} & \textbf{55} & \textbf{58} & \textbf{0.0006} & \textbf{0.002} \\
\textbf{X2} & \textbf{Uncertainty reported} & \textbf{63} & \textbf{80} & \textbf{85} & \textbf{95} & \textbf{100} & \textbf{88} & \textbf{0.0007} & \textbf{0.002} \\
Q8 & Statistical test applied & 6 & 5 & 10 & 10 & 25 & 12 & 0.13 & 0.13 \\
Q3 & All of the method's hyperparameters & 55\,{\scriptsize (47)} & 45 & 55 & 70 & 55 & 72 &  &  \\
Q4 & All of the baselines' hyperparameters & 35 & 16\,{\scriptsize (19)} & 45 & 25 & 40 & 29\,{\scriptsize (49)} &  &  \\
Q6 & Hyperparameter selection described & 26\,{\scriptsize (47)} & 30 & 40 & 40 & 35 & 48 &  &  \\
Q7 & Evaluation on held-out conditions & 33 & 25 & 20 & 20 & 55 & 45\,{\scriptsize (49)} &  &  \\
X4 & Baselines tuned with the same care & 13 & 26\,{\scriptsize (19)} & 35 & 15 & 25 & 16\,{\scriptsize (49)} &  &  \\
Y7 & Benchmark versions given & 23 & 25 & 20 & 40 & 20 & 26 &  &  \\
\midrule
\multicolumn{10}{@{}l}{\emph{Differences between periods (Newcombe, Fisher, Holm over the five tested criteria)}} \\
Code & Criterion & \multicolumn{3}{r}{\makecell[r]{2019-2020 against 2018\\{}gap [CI]}} & Holm $p$ & \multicolumn{3}{r}{\makecell[r]{2022-2024 against 2019-2020\\{}gap [CI]}} & Holm $p$ \\
\cmidrule(lr){1-10}
Q5 & Code linked & \multicolumn{3}{r}{+24 [+3,\,+42]} & 0.16 & \multicolumn{3}{r}{+17 [−3,\,+36]} & 0.49 \\
Y9 & Sensitivity to a hyperparameter & \multicolumn{3}{r}{−4 [−23,\,+16]} & 1 & \multicolumn{3}{r}{\textbf{+32 [+11,\,+50]}} & \textbf{0.03} \\
X7 & Compute quantified & \multicolumn{3}{r}{+8 [−11,\,+27]} & 1 & \multicolumn{3}{r}{+13 [−9,\,+32]} & 0.62 \\
X2 & Uncertainty reported & \multicolumn{3}{r}{+20 [+1,\,+37]} & 0.23 & \multicolumn{3}{r}{+15 [+1,\,+30]} & 0.23 \\
Q8 & Statistical test applied & \multicolumn{3}{r}{+1 [−10,\,+14]} & 1 & \multicolumn{3}{r}{+10 [−5,\,+25]} & 0.62 \\
\bottomrule
\end{tabularx}
\end{table}
